% Inactive reconstruction fragment; labels are prefixed pending integration.
% Chapter12 setting: compact connected X, Hermitian omega, 2<=m<n.

\begin{lemma}\label{draft631:proper-continuous-data}
Let $a>0$ and let smooth $H_j$ converge uniformly to a continuous
$H$. The smooth solutions of $H_m(w_j)=e^{aw_j+H_j}\omega^n$
converge uniformly to the unique continuous admissible solution of
$H_m(w)=e^{aw+H}\omega^n$. For two continuous data $H_1,H_2$,
the corresponding solutions satisfy
\[
 \|w_1-w_2\|_\infty\leq a^{-1}\|H_1-H_2\|_\infty.
\]
They also satisfy
$-\max_XH/a\leq w\leq-\min_XH/a$.
\end{lemma}
\begin{proof}
Smooth solutions exist by \cref{hessian:smooth-proper-existence}.
At a maximum of $w_1-w_2$, monotonicity of the smooth Hessian
operator gives $e^{aw_1+H_1}\leq e^{aw_2+H_2}$. Hence
$a(w_1-w_2)\leq H_2-H_1$ there. Exchanging the functions gives
the stated norm estimate for smooth data. Thus a uniformly
convergent sequence of smooth data gives a uniformly Cauchy
sequence of solutions. Its limit is admissible by the linear
distributional inequalities, and \cref{hessian:continuous-products}
passes the equation to the limit. The bounds follow from those
for the smooth solutions. Taking limits in the smooth norm estimate
proves it for continuous data. Finally,
\cref{hessian:hessian-exponential-order} in both directions gives
uniqueness among continuous admissible solutions.
\end{proof}

The proper equation can now be solved when its exponent depends
smoothly on the unknown potential, provided that constant barriers
bound the desired range and the dependence is strictly increasing.

\begin{theorem}\label{draft631:proper-variable-existence}
Let $A<B$ and let $f$ be smooth on a neighborhood of
$X\times[A,B]$. Suppose that, for every $x\in X$,
\[
 f(x,A)\leq0\leq f(x,B),\qquad
 0<\kappa\leq\partial_t f(x,t)\leq L\quad(A\leq t\leq B).
\]
There is a unique continuous admissible solution with values in
$[A,B]$ of
\[
 H_m(w)=e^{f(x,w)}\omega^n.
\]
This solution is smooth and strictly admissible.
\end{theorem}
\begin{proof}
Choose $a\geq L$ and let
$\mathcal C=\{v\in C(X):A\leq v\leq B\}$, a complete subset of
$C(X)$ in the uniform norm. For $v\in\mathcal C$ put
$H_v(x)=f(x,v(x))-av(x)$ and let $T(v)$ be the continuous proper
solution with exponent $aT(v)+H_v$, furnished by
\cref{draft631:proper-continuous-data}. The function $f(x,t)-at$
is nonincreasing in $t$, so
\[
 -aB\leq f(x,B)-aB\leq H_v(x)
                   \leq f(x,A)-aA\leq-aA.
\]
The zero-order bounds in that lemma give $A\leq T(v)\leq B$.
The mean-value theorem also gives
\[
 \|H_v-H_z\|_\infty\leq(a-\kappa)\|v-z\|_\infty,
 \qquad
 \|T(v)-T(z)\|_\infty
          \leq(1-\kappa/a)\|v-z\|_\infty.
\]
Since $0\leq1-\kappa/a<1$, the contraction theorem gives a unique
fixed point $w\in\mathcal C$. It is admissible because it equals
$T(w)$, and its equation is exactly the stated equation. Every
continuous admissible solution in this interval is a fixed point,
by uniqueness of the continuous proper equation with its prescribed
$H_w$. This proves uniqueness in the stated class.
\Cref{hessian:continuous-hessian-regularity} makes $w$ smooth.
Its form lies in the closed $m$-cone and has positive $m$th
symmetric polynomial; Maclaurin's inequalities then put it in the
open cone. Thus it is strictly admissible.
\end{proof}

The same strict increase controls solutions with a fixed Radon
right-hand measure. Here the bounded domination argument requires
the partial power identities on the background.

\begin{proposition}\label{draft631:proper-bounded-stability}
Assume $\ddc(\omega^{n-q})=0$ for $1\leq q\leq m$. Let $A<B$ and $f$ be
smooth near $X\times[A,B]$, with $\partial_t f\geq\kappa>0$
there, and let $\nu$ be a finite positive Radon measure.
Suppose bounded admissible $u,v$ have values in $[A,B]$ and satisfy
\[
 H_m(u)=e^{f(x,u)+h_1}\nu,\qquad
 H_m(v)=e^{f(x,v)+h_2}\nu,
\]
where $h_1,h_2$ are bounded Borel functions. Then
\[
 \|u-v\|_\infty
            \leq\kappa^{-1}\|h_1-h_2\|_{L^\infty(\nu)}.
\]
\end{proposition}
\begin{proof}
Put $K=\|h_1-h_2\|_{L^\infty(\nu)}$ and fix $s>K/\kappa$.
On $E_s=\{v+s<u\}$, the actual values $u,v$ lie in $[A,B]$,
so $f(x,u)-f(x,v)\geq\kappa s$. The equations and the essential
bound on $h_2-h_1$ therefore give, as restricted measures,
\[
 \mathbf1_{E_s}H_m(v+s)
 =\mathbf1_{E_s}H_m(v)
 \leq e^{K-\kappa s}\mathbf1_{E_s}H_m(u).
\]
The constant is smaller than one. Applying
\cref{hessian:strict-bounded-domination} to $v+s,u$ gives
$u\leq v+s$ pointwise. Let $s\downarrow K/\kappa$ and exchange
$u,v$ to obtain the norm bound. The exponent is evaluated only
at $u,v$; the shifted comparison potential need not take its
values in $[A,B]$.
\end{proof}

\begin{corollary}\label{draft631:proper-bounded-smoothness}
Assume $\ddc(\omega^{n-q})=0$ for $1\leq q\leq m$.
Under the hypotheses of \cref{draft631:proper-variable-existence},
every bounded admissible solution with values in $[A,B]$ is its
smooth solution. In particular, for every $a>0$ and smooth real
$H$, the equation $H_m(u)=e^{au+H}\omega^n$ has exactly one bounded
admissible solution, and that solution is smooth and strictly
admissible.
\end{corollary}
\begin{proof}
For a continuous admissible potential, smooth decreasing admissible
approximations converge uniformly on compact $X$ by Dini's theorem.
Thus \cref{hessian:joint-decreasing-products,hessian:continuous-products}
define the same Hessian measure on the continuous class.
The smooth solution of \cref{draft631:proper-variable-existence}
therefore solves the bounded equation as well. Apply
\cref{draft631:proper-bounded-stability} with $\nu=\omega^n$ and
$h_1=h_2=0$ to identify any other bounded solution in $[A,B]$
with it. For the final assertion, smooth existence is already
\cref{hessian:smooth-proper-existence}, and
\cref{hessian:proper-bounded-uniqueness} with
$\nu=e^H\omega^n$ gives uniqueness among all bounded solutions.
\end{proof}
